\section{Related Work}
\label{sec:related}

\subsection{Modeling Energetic Materials and Cocrystals}

Traditional modeling of energetic materials has relied on classical force fields and reactive potentials. Empirical force fields can capture crystalline structures and phonon spectra near equilibrium, but their fixed functional forms and limited parameterization often lead to poor transferability under strong compression or bond-breaking events. Reactive force fields introduce bond-order terms and empirical rules to allow reaction pathways \cite{van2001reaxff}, but they can be difficult to parameterize and may not faithfully reproduce \emph{ab initio} reaction energetics or charge transfer.

Energetic cocrystals introduce additional complexity: multiple molecular species coexist in the same lattice, with distinct electron affinities and ionization potentials. The cocrystal architecture enables tunable sensitivity and energy release by controlling intermolecular packing and charge transfer. Accurate modeling therefore requires a description that can resolve subtle differences in intermolecular interactions, follow bond-breaking and forming events, and adapt to evolving polarization and charge redistribution as shock loading proceeds \cite{valeeswaran2023charge,iosipescu2001experimental}. FPMD can in principle provide this detail, but remains limited in scale.

\subsection{Equivariant Neural Network Potentials}

Neural network interatomic potentials have evolved from early local descriptor models to modern architectures that tightly encode physical symmetries. Equivariant message passing networks such as NequIP, MACE, Allegro, and EquiformerV2 have demonstrated high accuracy across diverse materials settings \cite{batzner20223,batatia2022mace,kovacs2023mace,musaelian2023learning,liao2023equiformer}. Nonetheless, many such models still rely on finite interaction cutoffs, which can be a limiting approximation when long-range Coulomb interactions and configuration-dependent polarization remain important beyond the local neighborhood.

\subsection{Charge-Aware and Self-Consistent MLIPs}
\label{sec:qeq_review}

Charge equilibration (QEq) methods approximate partial charges by minimizing a quadratic functional derived from electronegativity equalization and hardness concepts \cite{rappe1991charge}. Classical and ML-augmented variants have long been used to introduce environment-dependent electrostatics \cite{gastegger2017machine,unke2019physnet,ko2021fourth}. More recent work has pushed charge-aware MLIPs significantly further, including equivariant local representations coupled to QEq, systematic studies of QEq-based ML potentials, and explicit long-range polarizable neural potentials.

For this manuscript, the key positioning point is the combination of design and validation choices that are particularly relevant for periodic energetic molecular crystals under shock-motivated distortion. First, CSC-MACE uses equivariant crystal-environment features to predict the diagonal QEq parameters inside the same architecture that produces the short-range energy, so the electrostatic subproblem is conditioned on the same local representation that governs the backbone potential instead of being introduced as a detached post hoc correction. Second, the charge solve is treated as part of the end-to-end periodic energy model: the same relaxed total energy is differentiated for energies, forces, and stresses, and charge neutrality is enforced over the full simulation cell rather than molecule by molecule. Third, the benchmark ladder explicitly separates three hypotheses that are often conflated in charge-aware MLIP papers: whether a short-range model is insufficient, whether a fixed-charge or static-QEq correction already captures most of the gain, and whether a fully self-consistent update changes the final target behavior enough to justify its extra complexity. The resulting contribution is to isolate what explicit self-consistency does and does not buy on a chemically demanding energetic-material target.
